\chapter{Proposed Architecture}
\label{ch:architecture}

This chapter presents the \padonap{} architecture, formalises its
threat model, and proves the safety invariants on which tier selection
rests. The architectural novelty is the closed loop between a multi-horizon
AI forecaster (S3) and an ONAP-compliant orchestration layer (S4), mediated
by a canonical JSON payload on DMaaP.

\section{System overview}
\label{sec:arch-overview}

\padonap{} is a four-stage pipeline (Figure~\ref{fig:arch}). Every 5-second
window, stage S1 scrapes per-device gNMI counters and (optionally) NetFlow
v5 records; stage S2 aggregates them into a 17-dimensional feature vector;
stage S3 runs XGBoost 7-class detection, Transformer+LSTM 4-horizon
forecasting, and emits a canonical \texttt{AIOutputPayload} JSON; stage S4
consumes this payload, evaluates the five-tier policy, allocates tenant
bandwidth under an LP SLA model, and instantiates VNFs through an ONAP SO
client.

\begin{figure}[H]
\centering
\begin{tikzpicture}[
  stage/.style={draw,rounded corners,minimum width=3cm,minimum height=1cm,align=center,font=\small},
  arr/.style={-Latex,thick},
]
\node[stage,fill=blue!10]   (s1) {S1 Telemetry\\(gNMI, NetFlow)};
\node[stage,fill=green!10,right=0.6cm of s1] (s2) {S2 Features\\(17 entropy/rate)};
\node[stage,fill=orange!15,right=0.6cm of s2] (s3) {S3 AI\\(XGB + Txf-LSTM)};
\node[stage,fill=red!12,right=0.6cm of s3] (s4) {S4 Orchestration\\(Policy+SLA+SO)};
\draw[arr] (s1)--(s2);
\draw[arr] (s2)--(s3);
\draw[arr] (s3)--(s4);
\node[below=0.5cm of s4,font=\small\itshape] {(closed loop: CLAMP/DMaaP feedback to S3)};
\end{tikzpicture}
\caption{\padonap{} four-stage pipeline. Each stage is independently replaceable.}
\label{fig:arch}
\end{figure}

\section{Threat model}
\label{sec:threat-model}

We formalise the adversary capability and the system's trust assumptions.

\begin{assumption}[Telemetry integrity]
\label{asm:telemetry}
The gNMI/NetFlow path from devices to stage S2 is authenticated and not
modifiable by the adversary. Compromise of telemetry itself is out of
scope; integrity can be provided by gNMI-over-TLS plus mTLS-authenticated
collectors.
\end{assumption}

\begin{assumption}[Control-plane trust]
\label{asm:control}
The ONAP SO/CLAMP/Policy plane and the DMaaP bus are operated in a trusted
administrative domain. Compromise of the orchestrator control plane is
out of scope.
\end{assumption}

\begin{definition}[Adversary capability]
\label{def:adv}
The adversary $\mathcal{A}$ controls a bounded-size botnet and can:
(a) emit volumetric floods (UDP, ICMP, amplification) up to link
saturation; (b) emit protocol attacks (SYN flood); (c) ramp attack rate
gradually to evade static thresholds; (d) launch low-and-slow attacks
staying below a specified rate. $\mathcal{A}$ cannot forge gNMI counters
nor insert messages on DMaaP.
\end{definition}

Out-of-scope threats include insider attacks, cryptographic attacks on
TLS, and compromise of the AI training data. The model in
Chapter~\ref{ch:ai} is trained offline on CICDDoS2019; data-poisoning
defenses are left for future work (see §\ref{sec:future}).

\section{Safety definition}
\label{sec:safety}

Proactive mitigation changes the agreement between defender and tenant:
a tier may be raised \emph{before} any attack has actually crossed a
threshold. We therefore require three properties.

\begin{definition}[Pre-positioning safety]
\label{def:safety}
For every tier transition $T_i \to T_j$ triggered proactively (i.e.~driven
by a forecast, not by a detection), the allocation to every benign tenant
$u$ must satisfy
\[
  a_u(T_j) \;\geq\; \text{floor}_u,
\]
where $\text{floor}_u$ is the SLA floor of $u$ and $a_u$ is its allocated
bandwidth.
\end{definition}

\begin{invariant}[URLLC floor]
\label{inv:urllc}
For every tier $T \in \{T_0,\dots,T_4\}$ and for every scenario,
$a_{\text{URLLC}}(T) \geq 150$ Mbps.
\end{invariant}

\begin{invariant}[Bounded rollback]
\label{inv:rollback}
After a proactive escalation $T_i \to T_j$ at window $w$, if the next
window $w+1$ reports tier $T_k < T_j$ for $k$ consecutive windows
(hysteresis $k{=}2$), the policy engine de-escalates within window $w+1+k$.
\end{invariant}

Invariants~\ref{inv:urllc} and~\ref{inv:rollback} are verified empirically
in Chapter~\ref{ch:eval} on all eight evaluation scenarios.

\section{Five-tier policy state machine}
\label{sec:tiers}

The mitigation action space is discretised into five tiers (Table~\ref{tab:tiers}).
Tier selection follows a deterministic map from the AI payload's
top-horizon probability and classification confidence, refined by a
hysteresis filter and a frequency guard to prevent flapping.

\begin{table}[H]
\centering
\small
\caption{Five-tier action ladder.}
\label{tab:tiers}
\begin{tabular}{@{}clll@{}}
\toprule
Tier & Label & Action & VNF profile \\
\midrule
$T_0$ & NORMAL   & no action                       & --- \\
$T_1$ & ALERT    & log + raise operator alert      & --- \\
$T_2$ & PREEMPT  & pre-position rate-limit VNF     & vnfd-ratelimiter-v1 \\
$T_3$ & MITIGATE & instantiate scrubber VNF, SFC   & vnfd-scrubber-v1 \\
$T_4$ & SEVERE   & blackhole + scrubber upstream   & vnfd-blackhole-v1 \\
\bottomrule
\end{tabular}
\end{table}

\paragraph{Hysteresis.}
Let $\tau(w)$ denote the tier issued at window $w$. The published tier is
the hysteresis-filtered value
\[
  \tilde\tau(w) \;=\; \max\!\left(\tau(w),\;\tilde\tau(w-1)\cdot\mathbb{1}[w - w_{\mathrm{last}} \leq k]\right),
\]
where $k{=}2$ is the hold window and $w_{\mathrm{last}}$ is the last window
in which the raw tier reached $\tilde\tau(w-1)$. This guarantees no rapid
oscillation between tiers and implements Invariant~\ref{inv:rollback}.

\section{LP multi-tenant SLA allocator}
\label{sec:sla-lp}

Let $n$ tenants share a total link capacity $C$. Denote by $d_u$ tenant
$u$'s demand, $f_u$ its floor, $w_u$ its priority weight, and $o(T)$ the
VNF overhead at tier $T$. We solve the LP
\begin{align}
  \max_{\mathbf{a}} &\quad \sum_{u=1}^n w_u \, a_u \label{eq:lp-obj}\\
  \text{s.t.} &\quad \sum_{u=1}^n a_u \;\leq\; C - o(T), \label{eq:lp-cap}\\
  &\quad f_u \;\leq\; a_u \;\leq\; d_u, \qquad \forall u. \label{eq:lp-bounds}
\end{align}
The primal is solved with HiGHS through SciPy; a proportional fallback
handles the case where scipy is unavailable. Constraint~\eqref{eq:lp-bounds}
directly enforces Invariant~\ref{inv:urllc} for $u = \text{URLLC}$.

\section{AIOutputPayload contract}
\label{sec:payload}

Stage S3 emits a canonical JSON object on DMaaP (topic
\texttt{PAD\_ONAP\_AI\_SIGNALS}):

\begin{lstlisting}[language=Python,caption={Skeleton of AIOutputPayload.}]
{
  "timestamp":     <epoch seconds>,
  "device_id":     "<str>",
  "detection":     {"attack_class": <int0..6>, "confidence": <float>},
  "forecast":      [{"horizon_s": 30, "p_attack": <float>}, ...],
  "top_features":  [{"feature": "<str>", "shap": <float>}, ...],
  "suggested_tier": <int0..4>,
  "proactive":     <bool>
}
\end{lstlisting}

Downstream ONAP consumers only need to parse this schema; the AI layer
can be replaced or upgraded without touching the policy engine. This
separation is what permits \padonap{} to claim independence from any
specific ML implementation.
